% CV for health-system and health-services research positions (e.g. delivery science fellowships).
% The general CV is main.tex; both share cv_preamble.tex and cv_header.tex.
\documentclass[10pt, letterpaper]{article}
% Shared preamble for main.tex (general CV) and cv_health_research.tex
% Packages:
\usepackage[
ignoreheadfoot, % set margins without considering header and footer
top=2 cm, % seperation between body and page edge from the top
bottom=2 cm, % seperation between body and page edge from the bottom
left=2 cm, % seperation between body and page edge from the left
right=2 cm, % seperation between body and page edge from the right
footskip=1.0 cm, % seperation between body and footer
% showframe % for debugging
]{geometry} % for adjusting page geometry
\usepackage{titlesec} % for customizing section titles
\usepackage{tabularx} % for making tables with fixed width columns
\usepackage{array} % tabularx requires this
\usepackage[dvipsnames]{xcolor} % for coloring text
\definecolor{primaryColor}{RGB}{0, 79, 144} % define primary color
\usepackage{enumitem} % for customizing lists
\usepackage{fontawesome5} % for using icons
\usepackage{amsmath} % for math
\usepackage[
pdftitle={Bipin Lekhak's CV},
pdfauthor={Bipin Lekhak},
pdfcreator={LaTeX with RenderCV},
colorlinks=true,
urlcolor=primaryColor
]{hyperref} % for links, metadata and bookmarks
\usepackage[pscoord]{eso-pic} % for floating text on the page
\usepackage{calc} % for calculating lengths
\usepackage{bookmark} % for bookmarks
\usepackage{lastpage} % for getting the total number of pages
\usepackage{changepage} % for one column entries (adjustwidth environment)
\usepackage{paracol} % for two and three column entries
\usepackage{ifthen} % for conditional statements
\usepackage{needspace} % for avoiding page brake right after the section title
\usepackage{iftex} % check if engine is pdflatex, xetex or luatex
\usepackage{etoolbox} % for \patchcmd and \AtBeginEnvironment

% Ensure that generate pdf is machine readable/ATS parsable:
\ifPDFTeX
    \input{glyphtounicode}
    \pdfgentounicode=1
    % \usepackage[T1]{fontenc} % this breaks sb2nov
    \usepackage[utf8]{inputenc}
    \usepackage{lmodern}
\fi



% Some settings:
\AtBeginEnvironment{adjustwidth}{\partopsep0pt} % remove space before adjustwidth environment
\pagestyle{empty} % no header or footer
\setcounter{secnumdepth}{0} % no section numbering
\setlength{\parindent}{0pt} % no indentation
\setlength{\topskip}{0pt} % no top skip
\setlength{\columnsep}{0cm} % set column seperation
\makeatletter
\let\ps@customFooterStyle\ps@plain % Copy the plain style to customFooterStyle
\patchcmd{\ps@customFooterStyle}{\thepage}{
\color{gray}\textit{\small Bipin Lekhak - Page \thepage{} of \pageref*{LastPage}}
}{}{} % replace number by desired string
\makeatother
\pagestyle{customFooterStyle}

\titleformat{\section}{\needspace{4\baselineskip}\bfseries\large\color{primaryColor}}{}{0pt}{}[\vspace{1pt}\titlerule]

\titlespacing{\section}{
% left space:
-1pt
}{
% top space:
0.3 cm
}{
% bottom space:
0.2 cm
} % section title spacing

\renewcommand\labelitemi{$\circ$} % custom bullet points
\newenvironment{highlights}{
\begin{itemize}[
    topsep=0.10 cm,
    parsep=0.10 cm,
    partopsep=0pt,
    itemsep=0pt,
    leftmargin=0.4 cm + 10pt
    ]
    }{
\end{itemize}
} % new environment for highlights

\newenvironment{highlightsforbulletentries}{
\begin{itemize}[
    topsep=0.10 cm,
    parsep=0.10 cm,
    partopsep=0pt,
    itemsep=0pt,
    leftmargin=10pt
    ]
    }{
\end{itemize}
} % new environment for highlights for bullet entries


\newenvironment{onecolentry}{
\begin{adjustwidth}{
    0.2 cm + 0.00001 cm
    }{
    0.2 cm + 0.00001 cm
    }
    }{
\end{adjustwidth}
} % new environment for one column entries

\newenvironment{twocolentry}[2][]{
\onecolentry
\def\secondColumn{#2}
\setcolumnwidth{\fill, 4.5 cm}
\begin{paracol}{2}
    }{
    \switchcolumn \raggedleft \secondColumn
\end{paracol}
\endonecolentry
} % new environment for two column entries

\newenvironment{header}{
\setlength{\topsep}{0pt}\par\kern\topsep\centering\linespread{1.5}
}{
\par\kern\topsep
} % new environment for the header

\newcommand{\placelastupdatedtext}{ % \placetextbox{<horizontal pos>}{<vertical pos>}{<stuff>}
\AddToShipoutPictureFG*{ % Add <stuff> to current page foreground
\put(
\LenToUnit{\paperwidth-2 cm-0.2 cm+0.05cm},
\LenToUnit{\paperheight-1.0 cm}
){\vtop{{\null}\makebox[0pt][c]{\small\color{gray}\textit{Last updated in October 2026}\hspace{\widthof{Last updated in October 2026}}
}}} %
} %
} %

% save the original href command in a new command:
\let\hrefWithoutArrow\href

% new command for external links:
\renewcommand{\href}[2]{\hrefWithoutArrow{#1}{\ifthenelse{\equal{#2}{}}{ }{#2 }\raisebox{.15ex}{\footnotesize \faExternalLink*}}}

% one line in the skills section: bold label, then the list
\newcommand{\skillline}[2]{\textbf{#1:} #2}




\begin{document}
    % Shared name and contact header, \input right after \begin{document}
    \newcommand{\AND}{\unskip
    \cleaders\copy\ANDbox\hskip\wd\ANDbox
    \ignorespaces
    }
    \newsavebox\ANDbox
    \sbox\ANDbox{}

    \placelastupdatedtext
    \begin{header}
        \textbf{\fontsize{24 pt}{24 pt}\selectfont Bipin Lekhak}

        \vspace{0.3 cm}

        \normalsize
        \mbox{{\color{black}\footnotesize\faMapMarker*}\hspace*{0.13cm}Rochester, NY} %
        \kern 0.25 cm %
        \AND %
        \kern 0.25 cm %
        \mbox{\hrefWithoutArrow{mailto:bl1282@rit.edu}{\color{black}{\footnotesize\faEnvelope[regular]}\hspace*{0.13cm}bl1282@rit.edu}} %
        \kern 0.25 cm %
        \AND %
        \kern 0.25 cm %
        \mbox{\hrefWithoutArrow{https://bipinlekhak.com.np/}{\color{black}{\footnotesize\faLink}\hspace*{0.13cm}bipinlekhak.com.np}}%

        \mbox{\hrefWithoutArrow{https://scholar.google.com/citations?user=z8T6GewAAAAJ}{\color{black}{\footnotesize\faGraduationCap}\hspace*{0.13cm}Google Scholar}}%
        \kern 0.25 cm%
        \AND%
        \kern 0.25 cm%
        \mbox{\hrefWithoutArrow{https://openreview.net/profile?id=~Bipin_Lekhak2}{\color{black}{\footnotesize\faFile*}\hspace*{0.13cm}OpenReview}}%
        \kern 0.25 cm%
        \AND%
        \kern 0.25 cm%
        \mbox{\hrefWithoutArrow{https://www.linkedin.com/in/bipin-lekhak/}{\color{black}{\footnotesize\faLinkedinIn}\hspace*{0.13cm}bipin-lekhak}}%
        \kern 0.25 cm%
        \AND%
        \kern 0.25 cm%
        \mbox{\hrefWithoutArrow{https://github.com/bipin-lekhak}{\color{black}{\footnotesize\faGithub}\hspace*{0.13cm}bipin-lekhak}} %
    \end{header}


    \vspace{0.3 cm - 0.3 cm}


    \section{Summary}
    \begin{onecolentry}
        PhD candidate developing \textbf{causal inference and machine learning methods for longitudinal health
        data}. My work estimates how a patient's outcomes would evolve under alternative treatment plans when
        treatments change over time and are confounded by the patient's condition, and adapts these predictions
        to each patient's own history. I have analysed real-world data from an NIH-funded study of in-home
        monitoring for heart-failure patients after hospital discharge, with clinical investigators at the
        University of Rochester Medical Center. Before the PhD, four years building machine learning systems in
        industry, including medical image analysis.
    \end{onecolentry}

    \vspace{0.2 cm}

    \begin{onecolentry}
        \textbf{Research interests:} causal inference with electronic health records, treatment-effect estimation
        under time-varying confounding, personalised prediction of treatment response, remote patient monitoring
        and care after hospital discharge, clinical risk prediction.
    \end{onecolentry}


    \section{Education}
    \begin{twocolentry}{
        \textit{Aug 2022 – Present}

        \textit{Expected June 2027}}
        \textbf{Rochester Institute of Technology}, Rochester, NY

        \textit{Ph.D. in Computing and Information Sciences}
    \end{twocolentry}

    \vspace{0.10 cm}
    \begin{onecolentry}
        \begin{highlights}
            \item Advisor: Prof. Linwei Wang, Computational Biomedicine Laboratory
            \item Dissertation proposal (defended July 2026): \textit{Learning Dynamics for Forecasting Future
            Trajectories under External Interventions}
        \end{highlights}
    \end{onecolentry}

    \vspace{0.2 cm}

    \begin{twocolentry}{
        \textit{Nov 2014 – Sep 2018}}
        \textbf{Pulchowk Campus, Institute of Engineering, Tribhuvan University}, Nepal

        \textit{B.E. in Electronics and Communication Engineering}
    \end{twocolentry}


    \section{Research Experience}

    \begin{twocolentry}{
        \textit{Rochester, NY}

        \textit{Aug 2022 – Present}}
        \textbf{Graduate Research Assistant}

        \textit{Computational Biomedicine Laboratory, RIT} (PI: Prof. Linwei Wang)
    \end{twocolentry}

    \vspace{0.10 cm}
    \begin{onecolentry}
        \begin{highlights}
            \item \textbf{Remote monitoring of heart-failure patients after discharge} (NIH/NINR R01NR018301, with
            the University of Rochester Medical Center). Analysed data from a sensor-equipped toilet seat that
            recorded cardiac measurements at home for 90 days after hospital discharge. Quantified enrollment and
            adherence, and identified why in-home deployments failed. Built a random-forest model that predicts
            hospitalisation within 21 days from 7-day windows of seat-derived measurements, including heart rate
            and heart-rate variability (AUC 0.74, sensitivity 57\%, specificity 75\%), and interactive dashboards
            for reviewing the study data. First author of the resulting Computing in Cardiology paper.
            \item \textbf{Treatment-response prediction under time-varying confounding (iASIDE).} Developed a model
            that predicts how a patient's outcomes will evolve under a planned sequence of treatments. It learns
            the course of the disease and the response to treatment as separate components, adapts both to the
            patient's history, and comes with an identifiability guarantee. Evaluated on MIMIC-III intensive-care
            records and a pharmacokinetic--pharmacodynamic simulator of tumour growth under chemotherapy and
            radiotherapy, with and without treatment-assignment bias.
            \item \textbf{Comparison with g-methods and their machine-learning extensions.} Built a reproducible
            benchmark of 10+ methods for time-varying treatments, including marginal structural models with
            inverse probability weighting, recurrent marginal structural networks, g-computation (G-Net) and
            balanced-representation methods (CRN, Causal Transformer), across prediction horizons and levels of
            confounding.
            \item \textbf{Heterogeneous treatment effects over time} (work in progress, with S.\ Poudel). Temporal
            X- and R-learners for estimating individual treatment effects from longitudinal data.
            \item \textbf{Reproducible research software.} Analysis pipelines in Python with experiment tracking
            (MLflow), configuration management and automated tests, run on high-performance computing clusters
            (SLURM).
        \end{highlights}
    \end{onecolentry}


    \section{Publications}

    \begin{onecolentry}
        \textit{Manuscripts under review}
    \end{onecolentry}

    \vspace{0.10 cm}

    \begin{samepage}
        \begin{twocolentry}{
            \textit{2026}
            }
            \textbf{iASIDE: Identifiable, Adaptable, and Separable Interventional Dynamics}

            \vspace{0.05 cm}

            \textbf{\textit{Bipin Lekhak}}, Yubo Ye, Suraj Poudel, Ruby Shrestha, Linwei Wang
        \end{twocolentry}
    \end{samepage}

    \vspace{0.3 cm}

    \begin{onecolentry}
        \textit{Peer-reviewed conference papers}
    \end{onecolentry}

    \vspace{0.10 cm}

    \begin{samepage}
        \begin{twocolentry}{
            \textbf{ICLR 2026}
            }
            \textbf{ProReGen: Progressive Residual Generation under Attribute Correlations}

            \vspace{0.05 cm}

            \mbox{Ruby Shrestha, Ajay Gopi, Casey Meisenzahl, \textbf{\textit{Bipin Lekhak}}, Linwei Wang}

            \vspace{0.05 cm}

            International Conference on Learning Representations.
            \href{https://openreview.net/forum?id=2LzYaW032Q}{OpenReview}
        \end{twocolentry}
    \end{samepage}

    \vspace{0.3 cm}

    \begin{samepage}
        \begin{twocolentry}{
            \textbf{CinC 2023}
            }
            \textbf{Novel In-Home Cardiac Monitoring for Heart Failure Patients}

            \vspace{0.05 cm}

            \textbf{\textit{Bipin Lekhak}}, Ryan Missel, Dillon Dzikowicz, Solomiya Rachynska, Wojciech Zareba, Linwei Wang

            \vspace{0.05 cm}

            Computing in Cardiology.
            \href{https://ieeexplore.ieee.org/abstract/document/10363960}{IEEE Xplore}
        \end{twocolentry}
    \end{samepage}


    \section{Industry Experience}

    \begin{twocolentry}{
        \textit{Kathmandu, Nepal}

        \textit{Sep 2018 – Jun 2022}}
        \textbf{Machine Learning Engineer}

        \textit{Fusemachines}
    \end{twocolentry}

    \vspace{0.10 cm}
    \begin{onecolentry}
        \begin{highlights}
            \item Built deep-learning segmentation of medical images (U-Net, Mask R-CNN, Detectron).
            \item Reconstructed 3D anatomical structures from 2D images by optimisation over statistical shape
            models.
            \item Forecast product sales from historical data with multi-stage statistical and machine learning
            models (LSTMs, Prophet, gradient-boosted trees, Bayesian models in PyMC).
        \end{highlights}
    \end{onecolentry}


    \section{Teaching}

    \begin{twocolentry}{
        \textit{Fall 2025}}
        \textbf{Teaching Assistant}, GCCIS-820 Quantitative Foundations, RIT

        \textit{PhD core course: linear algebra, optimisation and statistical inference}
    \end{twocolentry}

    \vspace{0.10 cm}
    \begin{onecolentry}
        \begin{highlights}
            \item Supervised and graded team projects and reports; built the tool used to form project teams.
        \end{highlights}
    \end{onecolentry}


    \section{Skills}

    \begin{onecolentry}
        \skillline{Causal inference}{potential outcomes, time-varying confounding, inverse probability weighting
        and marginal structural models, g-computation, meta-learners (X- and R-learners)}
    \end{onecolentry}

    \vspace{0.10 cm}

    \begin{onecolentry}
        \skillline{Statistical and ML modelling}{longitudinal and time-series models, state-space and
        latent-variable models, Bayesian modelling, clinical risk prediction (random forests, gradient boosting),
        deep learning}
    \end{onecolentry}

    \vspace{0.10 cm}

    \begin{onecolentry}
        \skillline{Health data}{MIMIC-III electronic health records, in-home remote-monitoring sensor data,
        medical images}
    \end{onecolentry}

    \vspace{0.10 cm}

    \begin{onecolentry}
        \skillline{Software}{Python (pandas, scikit-learn, PyTorch, PyMC), R, SQL, MATLAB, LaTeX; Git, MLflow,
        SLURM and HPC clusters}
    \end{onecolentry}

\end{document}
